\chapter{Estrutura de pastas e reprodutibilidade}
\label{app:pastas}

Este apêndice descreve a organização do repositório da dissertação e a ordem
em que os scripts reproduzem os resultados. As definições das variáveis estão
no Apêndice~\ref{app:variaveis}, e os trechos de código, no
Apêndice~\ref{app:scripts}.

\section{Estrutura de diretórios}

\begin{verbatim}
bitcoin-volatility-risk-premium/
├── data/                      # séries e conjuntos de dados
│   ├── btc_prices.csv         # preço de fechamento diário (Binance)
│   ├── dvol_30d_full.csv      # índice DVOL diário (Deribit)
│   ├── vrp_30d_dataset.csv    # VH, IV, proxy e BVRP
│   ├── vrp_with_targets.csv   # + retornos futuros: amostra completa
│   ├── ml_dataset_T4.csv      # alvo e 17 variáveis: amostra de modelagem
│   ├── regimes_T9.csv         # cortes e regimes de volatilidade
│   ├── previsoes_bvrp_T5.csv  # BVRP previsto pelos cinco modelos
│   └── previsoes_bvrp_T11.csv # BVRP previsto pelo modelo em árvore
├── code/                      # coleta, construção, análises e geradores
│   ├── *_utils.py             # divisão temporal, Newey-West e avaliação
│   ├── publicar_*.py          # geradores das tabelas e figuras
│   └── test_*.py              # testes
├── scripts/                   # teste de H1 e conferência do LaTeX
├── outputs/                   # saídas de cada análise, por script
│   └── T1/ T2/ T4/ T5/ T9/ T10/ T11/ T12/ C3/ C7/ A1/ MCS/
├── figs/                      # figuras usadas no texto, por tema
├── tables/                    # tabelas usadas no texto, por tema
│   └── dados/ previsao_bvrp/ retornos_linear/
│       retornos_nao_linear/ estrategias/
├── chapters/                  # capítulos, nota metodológica e apêndices
├── frontmatter/               # documento principal, preâmbulo e .bib
├── docs/                      # registro das decisões metodológicas
└── archive/                   # arquivos fora do documento
\end{verbatim}

As pastas \arq{figs/} e \arq{tables/} têm os mesmos cinco temas: dados
(Capítulo~\ref{chap:dados}), previsão do BVRP
(Capítulo~\ref{chap:previsao_bvrp}), retornos com o modelo linear
(Capítulo~\ref{chap:retornos_linear}) e com o modelo em árvore
(Capítulo~\ref{chap:retornos_nao_linear}) e estratégias
(Apêndice~\ref{app:estrategias}). Os demais scripts de \arq{code/} e
\arq{scripts/} não produzem resultados usados no texto.

\section{Ordem de execução}

\begin{enumerate}
    \item \textbf{Coleta}: \arq{update_btc_prices.py} e
          \arq{download_dvol_deribit.py} gravam
          \arq{data/btc_prices.csv} e \arq{data/dvol_30d_full.csv};
          \arq{fix_btc_gap_2023_03.py} \arq{--write} completa o mês de março
          de 2023 na série de preços.
    \item \textbf{Construção}: \arq{build_vrp_dataset.py},
          \arq{build_targets.py}, \arq{build_ml_dataset_T4.py} e
          \arq{regimes_T9.py}, nessa ordem. O primeiro interrompe a
          execução se faltar alguma data nas séries diárias.
    \item \textbf{Dados e existência do prêmio}:
          \arq{build_desc_stats_T1.py}, \arq{plot_cap3_T1.py},
          \arq{regimes_T9.py} e \arq{plot_vrp_vs_return_T12.py},
          com a opção \arq{--publicar-cap3}, que copia as tabelas e as
          figuras para \arq{tables/dados/} e \arq{figs/dados/};
          \arq{scripts/test_h1_bvrp_mean.py} e
          \arq{descritivas_C3_regimes.py}.
    \item \textbf{Previsão do BVRP}: \arq{previsao_bvrp_T5.py} e
          \arq{interacao_regime_T9.py}; depois, \arq{mcs_previsao_T5.py} e
          \arq{publicar_cap5.py}.
    \item \textbf{Retornos}: \arq{retorno_bvrp_T10.py},
          \arq{retorno_bvrp_T11.py} e \arq{descritivas_C7.py}; depois,
          \arq{publicar_cap6.py} e \arq{publicar_cap7.py}. O
          \arq{retorno_bvrp_T11.py} vem antes do
          \arq{publicar_cap6.py} porque calcula também o teste conjunto
          do modelo linear.
    \item \textbf{Estratégias}: \arq{estrategias_A1.py}; depois,
          \arq{publicar_apendice.py}.
    \item \textbf{Conferência}: os testes \arq{code/test_*.py} e
          \arq{scripts/verificar_latex.py}, que confere os arquivos
          incluídos, as referências cruzadas e as citações do documento.
\end{enumerate}

Salvo indicação, os scripts estão em \arq{code/} e rodam a partir da raiz
do repositório, com \texttt{python code/<script>.py}.

\section{Reprodutibilidade}

Os modelos com componente aleatório usam semente fixa: 42 nas árvores e nas
permutações do placebo, e 20261001 nas reamostragens do bootstrap, as mesmas
nos Capítulos~\ref{chap:retornos_linear} e~\ref{chap:retornos_nao_linear}, e nas
do \textit{model confidence set}. As
tabelas e as figuras do texto são gravadas pelos scripts, sem edição manual.
A Nota Metodológica resume as escolhas de implementação.
